% s-239-p239-m001; original PDF 239
\Aphorism{s-239-p239-m001}{331}
Strenge Dichtformen, wie das Sonett, thun ganz gewiss der sachlichen Feinheit Schaden; Reim und Rhythmus zwingen, in der Mehrzahl der Fälle, einiges mehr, einiges weniger, vieles anders zu sagen als in der ersten Conception lag. Die Chancen des Gelingens verringern sich, in Folge der „Abtrift“: es ist unwahrscheinlich, dass der mitzutheilende Inhalt innerhalb der willkürlich gewählten Form wirklich mittheilbar werde, dass er sie dulde, ja gar zu fordern scheine — und dies erst wäre erreichte Vollkommenheit. — Aber eben in dieser Gefahr der strengen und absoluten Formen liegt ihr disciplinarischer Werth. Es muss vor allem verhindert werden, dass dem Dichter das Versemachen leicht falle, dass er sich „ohne Noth“ dichten lasse; dazu eben dient irgend ein erschwerender Formalismus, der, sobald er erlernt ist und nicht mehr straff genug anzieht, mit einem härteren zu vertauschen ist. Bei dieser Schaffensart wird Vieles als blosser Versuch unter den Tisch fallen und, wie gesagt, weniger gelingen, als wenn allein der sachlichen Ausarbeitung Recht widerführe; aber eben diese Auslese thut Noth. Bisher vollzog sie der Leser, der Herausgeber, bestenfalls der Dichter selbst; nun wird sie bereits in der dichterischen Methode vollzogen werden, und was die Werkstatt verlässt, wird ganz andere Bürgschaften für Ächtheit und Dauer bieten, als zuvor auch nur verlangt wurden.
\par

